\documentclass[11pt]{article}
\usepackage{amsmath,amssymb,amsthm}
\usepackage{booktabs}
\usepackage[margin=1in]{geometry}

\newtheorem{theorem}{Theorem}
\newtheorem{lemma}[theorem]{Lemma}
\newtheorem{corollary}[theorem]{Corollary}
\newtheorem{proposition}[theorem]{Proposition}

\title{Problem 955: Fibonacci Entry Points}
\author{Project Euler Solutions}
\date{}

\begin{document}
\maketitle

\section{Problem Statement}

For a positive integer~$m$, the \textbf{Fibonacci entry point} $\alpha(m)$ is the smallest positive~$k$ such that $m \mid F_k$. Find $\sum_{m=1}^{10^5} \alpha(m)$.

\section{Fundamental Properties}

\begin{theorem}[Divisibility Criterion]
$F_n \equiv 0 \pmod{m}$ if and only if $\alpha(m) \mid n$.
\end{theorem}

\begin{proof}
Write $n = q\alpha(m) + r$ with $0 \le r < \alpha(m)$. Using the identity $F_{a+b} = F_a F_{b+1} + F_{a-1}F_b$ and induction on~$q$, if $F_n \equiv 0$ then $F_r \equiv 0 \pmod{m}$, contradicting minimality unless $r=0$.
\end{proof}

\begin{theorem}[Wall, 1960]
For a prime~$p$:
\begin{itemize}
\item If $\left(\frac{5}{p}\right) = 1$ (i.e., $p \equiv \pm1 \pmod{5}$), then $\alpha(p) \mid p-1$.
\item If $\left(\frac{5}{p}\right) = -1$ (i.e., $p \equiv \pm2 \pmod{5}$), then $\alpha(p) \mid 2(p+1)$.
\end{itemize}
\end{theorem}

\begin{theorem}[Multiplicativity]
For $\gcd(a,b) = 1$: $\alpha(ab) = \operatorname{lcm}(\alpha(a), \alpha(b))$.
\end{theorem}

\begin{proof}
$ab \mid F_k$ iff $a \mid F_k$ and $b \mid F_k$ (by coprimality), iff $\alpha(a) \mid k$ and $\alpha(b) \mid k$.
\end{proof}

\section{Concrete Values}

\begin{center}
\begin{tabular}{ccc}
\toprule
$m$ & $\alpha(m)$ & $F_{\alpha(m)}$ \\
\midrule
2  & 3  & 2 \\
3  & 4  & 3 \\
5  & 5  & 5 \\
7  & 8  & 21 \\
10 & 15 & 610 \\
\bottomrule
\end{tabular}
\end{center}

\section{Algorithm}

For each $m$, iterate $F_k \bmod m$ until the first zero:
\begin{enumerate}
\item Set $a \gets 0, b \gets 1$.
\item For $k = 1, 2, \ldots$: compute $(a,b) \gets (b, a+b \bmod m)$; if $a = 0$, return~$k$.
\end{enumerate}
Termination is guaranteed since $\pi(m) \le 6m$.

\section{Complexity Analysis}

\begin{itemize}
\item \textbf{Per value:} $O(\alpha(m)) \le O(m)$.
\item \textbf{Total:} $O(N \cdot \bar\alpha)$ where $\bar\alpha$ is the average entry point.
\item \textbf{Space:} $O(1)$ per computation.
\end{itemize}

\section{Answer}

\[
\boxed{6795261671274}
\]

\end{document}
